% Generated draft: requires mathematical and source review.
\documentclass[11pt]{article}
\usepackage[margin=25mm]{geometry}
\usepackage{fontspec}
\setmainfont{Latin Modern Roman}
\usepackage{amsmath,amssymb,mathtools,unicode-math}
\setmathfont{Latin Modern Math}
\usepackage{enumitem,needspace,xurl,hyperref}
\setlength{\parindent}{0pt}
\setlength{\parskip}{5pt}
\setlength{\emergencystretch}{3em}
\newcommand{\N}{\mathbb{N}}
\newcommand{\Z}{\mathbb{Z}}
\newcommand{\Q}{\mathbb{Q}}
\newcommand{\R}{\mathbb{R}}
\newcommand{\C}{\mathbb{C}}
\newcommand{\F}{\mathbb{F}}
\newcommand{\A}{\mathbb{A}}
\newcommand{\PP}{\mathbb{P}}
\newcommand{\E}{\mathbb{E}}
\newcommand{\eps}{\varepsilon}
\newcommand{\dd}{\,\mathrm{d}}
\begin{document}
\section*{UnsolvedMath ID: 30004972}
\textbf{Source problem code:} catalogue entry
\subsection{Definitions and mathematical statement}
Let \(n\geq 1\) be an integer and write \([n]=\{1,\ldots,n\}\). A Latin square of order \(n\) is a function \(L:[n]\times[n]\to[n]\) such that each symbol occurs exactly once in each row and exactly once in each column. A partial transversal is a set \(T\subseteq[n]\times[n]\) such that, for distinct \((i,j),(i',j')\in T\), one has \(i\neq i'\), \(j\neq j'\), and \(L(i,j)\neq L(i',j')\). Its length is \(|T|\). A full transversal has length \(n\).

\textbf{Near-transversal conjecture (Ryser--Brualdi--Stein, as attributed in the supplied record).} For every integer \(n\geq1\) and every Latin square \(L\) of order \(n\), there exists a partial transversal \(T\) satisfying
\[
|T|=n-1.
\]
For \(n=1\), the empty set is permitted. Equivalently, every Latin square of order \(n\) has a partial transversal of length at least \(n-1\), since entries can be discarded from a longer partial transversal.

\textbf{Source qualifications.} The original excerpt uses the phrase ``a transversal which uses all but at most one symbol''; the supplied clean statement interprets this as a partial transversal of length at least \(n-1\). The dated literature assessment also discusses a full-transversal assertion for odd orders. That assertion would require \(|T|=n\) whenever \(n\) is odd, and is a separate, stronger requirement than the near-transversal statement formulated here. The record labels the problem partially solved and labels its statement \texttt{corrected\_verified}; these are supplied curation labels, not independent verification in this draft.
\subsection{Short English statement}
Does every Latin square of positive order \(n\) contain \(n-1\) entries in distinct rows and distinct columns, all carrying different symbols? The supplied dated assessment reports results for sufficiently large orders; this draft preserves the universal conjecture and requires human mathematical review.
\subsection{Sources}
\begin{itemize}
\item Local UnsolvedMath registry, \texttt{data/unsolvedmath/problems.json}, permanent ID \(30004972\), source code \texttt{catalogue entry}, title: Ryser--Brualdi--Stein Near-Transversal Conjecture. The supplied record belongs to the Oberwolfach Reports Open Problems collection and lists proposed year \(2022\), with no individual proposer specified.
\item Supplied report citation: \emph{Graph Theory (2023)}, Oberwolfach Reports, DOI \texttt{10.4171/owr/2022/1}, pp.\ 5--78. \url{https://doi.org/10.4171/owr/2022/1}. The literature-source label calls this \emph{Oberwolfach Report 1/2022}. The discrepancy between the supplied report year and the DOI/report label remains unresolved.
\item The supplied original excerpt attributes the conjecture to Ryser, Brualdi and Stein, citing report references [10, 5, 11]. It appears alongside the contribution \emph{Short proofs of rainbow matching results}, attributed to Benny Sudakov, joint work with David Munh\'{a} Correia and Alexey Pokrovskiy. The excerpt also mentions an Aharoni--Berger conjecture with reference [1], but supplies no statement of that separate conjecture. Bibliographic entries for these numbered references were not supplied.
\item Supplied literature source: ``Montgomery, large-even-order Ryser--Brualdi--Stein'', \url{https://arxiv.org/abs/2310.19779}.
\item Literature checked at: \texttt{2026-08-22}. This is supplied dated curation, not independently verified. Its assessment states that the near-transversal part is proved for all sufficiently large orders, including large even \(n\), while the full-transversal assertion for odd orders and finite residual cases remain open. The background describes this assessment as dated literature triage derived from the supplied Oberwolfach curation and requests independent verification of time-sensitive status and cited claims.
\item The supplied statement-verification field says the statement was checked against the official Oberwolfach report and preserves its hypotheses, notation and open target. This draft records that provenance claim without claiming to have visited or checked the report.
\end{itemize}
\end{document}
